\section{Selected Job Position and Job Description}
\newcommand{\jobad}{google-about-careers-applications-jobs-results-121141635160384198-software-engineer-android-sre}

\paragraph{The role.}
Google's SRE book says that SRE teams hire software engineers to run products and to build systems that do work sysadmins would otherwise do, often manually \cite{sre-book}. The SRE Workbook says the basic tenet of SRE is that doing operations well is a software problem, and describes SRE as both a job role and a set of practices \cite{workbook}. The SRE book also says SRE can be viewed as a specific implementation of DevOps with some idiosyncratic extensions \cite{sre-book}. A Google job advertisement from one team, Software Engineer, Android SRE, says that SRE combines software and systems engineering to build and run large-scale, massively distributed, fault-tolerant systems \cite{\jobad}.

\paragraph{Duties.}
An SRE team is generally responsible for the availability, latency, performance, efficiency, change management, monitoring, emergency response and capacity planning of its services \cite{sre-book,workbook}. Google caps operational work (toil) at 50\% of an SRE's time, and the rest should go to project work that uses coding skills \cite{sre-book}. The job advertisement lists writing product or system development code, triaging and debugging issues by analysing their impact on hardware, network or service operations, and taking part in design reviews \cite{\jobad}.

\paragraph{Daily and weekly tasks.}
Operational load falls into three categories: pages, tickets and ongoing operational activities \cite{sre-book}. Being on call is one of the defining characteristics of SRE at Google, and during a shift SREs diagnose, mitigate, fix or escalate incidents and handle nonurgent production duties \cite{workbook}. After a page is received and acknowledged, the on-call engineer is expected to triage the problem and work toward its resolution, escalating as needed \cite{sre-book}. Typical paging response times are 5 minutes for user-facing or time-critical services and 30 minutes for less time-sensitive systems \cite{sre-book}. SRE teams must write postmortems with a full timeline after significant incidents \cite{sre-book}. The books describe SRE teams in general. They give no day-by-day schedule for one SRE, and we do not infer one.

\paragraph{Projects.}
The books describe project work in four groups. The first is internal tools and automation, such as binary rollout mechanisms and monitoring \cite{sre-book}. The second is capacity planning, where Auxon automates the work for services in Google production \cite{sre-book}. The third is onboarding services: the Production Readiness Review (PRR) is a prerequisite for an SRE team to accept responsibility for managing a service in production \cite{sre-book}. The fourth is launch support, where in one example a special team of Launch Coordination Engineers coordinated the groups involved \cite{sre-book}. A Workbook case study also describes automating the repair of failed datacenter line cards \cite{workbook}. We found no source on projects given to student SREs.

\paragraph{Education and qualifications.}
The SRE book says 50--60\% of SREs are hired through the standard procedure for Google Software Engineers. The other 40--50\% are very close to those qualifications and also have technical skills that are useful to SRE but rare among software engineers, most commonly UNIX system internals and networking (Layer 1 to Layer 3) \cite{sre-book}. The job advertisement, which is not for an internship, lists as minimum qualifications a Bachelor's degree in Computer Science, a related field or equivalent practical experience, and 2 years of software development experience in one or more programming languages. It prefers a Master's degree and 2 years of experience troubleshooting large-scale distributed systems \cite{\jobad}. We found no source on certificates, a portfolio or open-source experience that would help an applicant.